\documentclass[12pt,a4paper]{report}
\usepackage{charte_graphique_INSA/classeRapport} % template INSA
\usepackage{float}
\usepackage{minted}
\usepackage{amssymb}
\usepackage{pifont}
\usepackage{appendix}
\usepackage[backend=biber]{biblatex}
\usepackage{csquotes}
\usepackage{longtable}
\setlongtables
\addbibresource{rapport.bib}


\hypersetup
{
 %%   Champs dans les propriétés du document
 pdftitle={Rapport de Stage Ingénieur -- Prénom NOM -- 20xx}, % titre dans les propriétés du pdf
 pdfauthor={Prénom NOM}, % auteur dans les propriété du pdf
 pdfsubject={Rapport du Stage Ingénieur de Prénom NOM chez XXXXX en 20xx}, % sujet dans les propriétés du pdf
 pdfkeywords={}, % mots-clés dans les propriétés du pdf
}

\begin{document}

\PageDeGarde%
{logo_entreprise}% Logo entreprise
{Titre de votre document}% Titre du document
{Sous Titre}% Sous Titre
{Auteurs}% Auteurs
{Text en bas à droite}% Texte en bas à droite

\Page{INSALogo}{logo_entreprise}

\newpage
\thispagestyle{empty}
~ %Page de garde blanche et vierge (pour les notes du lecteur) après la page de couverture
\newpage 

\chapter*{Remerciements}
\setcounter{page}{1}
\addcontentsline{toc}{chapter}{Remerciements}
\chaptermark{Remerciements}
\input{remerciements} %1 page


\chapter*{Lexique}
\addcontentsline{toc}{chapter}{Lexique}
\chaptermark{Lexique}
\input{lexique} %1 page

% Sommaire avec pagination
\setcounter{tocdepth}{2}
\tableofcontents %2 pages

% contexte du stage (lieu, service, présentation rapide de l'organisme) et objectifs du stage
\chapter*{Introduction}
\addcontentsline{toc}{chapter}{Introduction}
\chaptermark{Introduction}
\input{Introduction}

\section*{Présentation de l'entreprise}
\addcontentsline{toc}{section}{Présentation de l'entreprise}
\sectionmark{Présentation de l'entreprise}
\input{presentation_entreprise}

\section*{Présentation du sujet du stage}
\addcontentsline{toc}{section}{Présentation du sujet du stage}
\sectionmark{Présentation du sujet du stage}
\input{sujet_stage}

% Le travail effectué [30-45 pages]
% Vous devez montrer les points conceptuels et techniques importants :
%     Étude du cahier des charges
%     Propositions et critiques de solutions
%     Description complète de la solution choisie
%     Mise en œuvre
%     Résultats obtenus (critique)

% 

\chapter{Nom de la première partie}
\input{part_1}

\chapter{Nom de la deuxième partie}
\input{part_2}

\chapter{Nom de la troisième partie}
\input{part_3}

\chapter{Impact environnemental et sociétal}
\input{ddrs}

\chapter*{Conclusion}
\addcontentsline{toc}{chapter}{Conclusion}
\chaptermark{Conclusion}
\input{conclusion} %2 pages


% \chaptermark{Bibliographie}
% Bibliographie [1-2 pages]
\printbibliography %1 page
\addcontentsline{toc}{chapter}{Bibliographie}

\begin{appendix}
\input{annexes} %2 pages
\end{appendix}

\input{resume} %1 page

\end{document}